% included by aGMX_clustering.tex: the changes of 25-26 October 2026 (phase 1, physics audit, few-electron detection,
% the few-electron candidate catalogue). Numbers from the validation runs named in each paragraph.
%======================================================================================================
\section{Faster EM, few-electron finder and fairer tests (phase 1, 25 October)}\label{sec:phase1}

\paragraph{Compiled EM.} The EM iteration of section~\ref{sec:heart} is fused into one compiled pass (numba): the
model $M$ is accumulated first, then the responsibility-weighted moments, with the same footprints and the same
M-step. The results are identical to $10^{-10}$\,keV on a clean and on the crowded CCD; the crowded CCD takes 15\,s
instead of 32\,s. \code{AGMX\_NUMBA=0} selects the numpy version.

\paragraph{Few-electron finder with the full pixel likelihood.} A group of touching pixels with $n_i\ge1$ counted
electrons is scored by
\begin{equation}
T=-\sum_{i\in g}\ln P(n_i\,|\,\lambda_i,s_i)-\ln N_\mathrm{shapes}(m),\qquad
P(n\,|\,\lambda,s)=\sum_k \mathrm{Pois}(k;\lambda)\Big[\Phi\big(\tfrac{n+\frac12-k}{s}\big)-\Phi\big(\tfrac{n-\frac12-k}{s}\big)\Big],
\end{equation}
the probability of the observed rounded counts under dark charge plus read noise (rounding miscounts included), with
the number of 8-connected shapes of $m$ pixels as the trial factor of the shape. The threshold is the value exceeded
by the wanted number of fake groups ($0.08$ per Mpix by default, the earlier operating point) in blank images
simulated from the image's own background; the analytic bound $\ln(N_\mathrm{pix}/F)$ is used only when the blanks
have too few groups. Section~\ref{sec:lowdark} gives the performance at the lowest dark currents.

\paragraph{Injection tests: objects away from real events.} The injected objects were placed where their centre was
quiet; a 15--35-px track could still cross a real event, and every method then ``failed'' on it. Objects are now
placed only where their whole footprint (+2\,px) avoids real events (\code{INJECT\_FOOTPRINT=0} for the earlier
placement, used for the wide test so that all methods stay comparable). On the injected tracks of the real Am-241
images (EXT3, the denser CCD) the muon success of every method goes from $\approx0.2$ to $1.00$: the earlier numbers
measured overlap, not clustering. With the new placement, \AGMX{} on EXT1: muon $0.93\to0.99$, worm $0.93\to0.99$,
hot column $0.74\to0.90$, pairs 3/4/6\,px $0.65/0.90/0.94$ (new objects); on the dense EXT3 not every object kind finds
a clean place any more (alphas, hot columns, low-E). On the 878-image wide test phase~1 changes no X-ray result (real
K$\alpha$ identical to 0.0\,eV on every image-CCD, injections within 0.006).

%======================================================================================================
\section{Physics audit (25--26 October)}\label{sec:audit}

Every step of \AGMX{} and of the other methods was checked for physical and statistical soundness (simulations of
skipper images with exactly known electrons; code review). The findings and what was done:

\begin{table}[H]\centering\footnotesize
\begin{adjustbox}{max width=\linewidth}
\begin{tabular}{p{0.6cm}p{6.3cm}p{6.4cm}p{3.2cm}}
\toprule
id & finding & fix & effect (simulation) \\
\midrule
A1 & a row counts as ``comb resolved'' when $|\langle e^{2\pi iQ}\rangle|>0.018$, below the noise floor $\sqrt{\pi/4N}$ of
pure noise: Gaussian noise of 1.5\,\el{} was fitted as $s\approx0.4$\,\el, the switch did not fire &
threshold $\max(0.018,4/\sqrt N)$; rows whose comb is washed out by event charge take $o,\lambda,s$ from the nearest
resolved rows (the read noise belongs to the amplifier); the pooled broad-noise fit only when $<5\,\%$ of the rows
are resolved & fakes on Gaussian noise 7966 $\to$ 1 per Mpix; crowded CCD unchanged \\
I1 & the read-noise estimate (MAD of $|Q|<0.5$) cannot exceed $1.4826\times0.25=0.37$\,\el &
robust spread of all active pixels when the first estimate is $>0.3$ & 0.37 $\to$ 1.50 / 3.00\,\el{} \\
A2 & rounding counts noise $>+0.5$\,\el{} as an electron: $\lambda$ high by $p=\Phi^c(0.5/s)$ (and the row mean, so
every energy $-N_\mathrm{pix}p$) & $\lambda\leftarrow(\lambda-p)/(1-2p)$ on comb-fitted rows (offset unbiased) &
$s=0.3$, $\lambda=0.001$: 0.051 $\to$ 0.004 \\
--- & $\lambda$ clipped at $10^{-3}$\,\el/pix (above DAMIC-M dark currents) & floor $10^{-6}$; detection uses the
measured single-electron occupancy (section~\ref{sec:lowdark}) & \\
B1 & bright hot columns hidden by the cluster mask (not flagged); flagged ones only had the mean subtracted (no
variance) & column excess (clipped mean, no mask) added to the column's $\lambda$, so mean, variance and the Poisson
tests follow & +2\,\el/pix column: 50 $\to$ 0 fake few-electron clusters \\
C2 & counting mode for $s<0.5$\,\el{} although rounding miscounts are $\sim10\,\%$ at $s=0.4$ & counting only for
$s<0.35$\,\el & $s=0.45$ blank: 313 $\to$ 0 fakes per Mpix \\
D1 & the pixel model samples a Gaussian of variance $\sigma^2+\frac1{12}$ at the pixel centre; its sum is 1.38 / 0.68
for $\sigma=0.2$ (centred / half-pixel) & exact pixel integrals: phase~2 (section~\ref{sec:phase2}) & sharp X-rays fitted with one cloud \\
D2 & the E-step floor $\varepsilon=\sqrt V$ pulls the fitted width low (1--4\,\% for Fe-55, up to 18\,\% for
200\,\el{} at 1\,\el/pix): the width is the depth estimator & reported width = second moment of the cluster's shared
charge (\code{sxy}); the EM width is kept as \code{sxy\_fit} & Fe-55: 0.2--1.8\,\% instead of $-1$ to $-4.5$\,\% \\
G1 & muon label from geometry only & a straight track is a muon only above 2.5\,keV per projected pixel (a minimum
ionising particle loses 0.27\,keV/$\mu$m, $\ge4$\,keV per 15-$\mu$m pixel) & physical muons 33/40 muon; the old
too-faint test muons $\to$ ``extended'' \\
G2 & saturated alphas not flagged & flag \code{saturated} (pixels at the image's maximum plateau) & \\
\bottomrule
\end{tabular}
\end{adjustbox}
\caption*{Physics-audit findings and fixes (\code{AGMX\_PHYS\_FIXES=0} restores the earlier behaviour).}
\end{table}

\paragraph{Real data.} 104 image-CCDs (every 12th Fe-55 image-CCD, every 6th Am-241 and all PD06/DM12), fixes on vs
off with the same injected objects: Fe-55 data unchanged (injections equal to 0.004; front-side K$\alpha$ $-0.2$\,eV
median, at most 7\,eV). The noisy PD06/DM12 images (read noise $\approx3.5$--4\,\el, now measured correctly) keep the
WADERs-like branch with the earlier thresholds (pixel $>4.5$\,\el, seed $\approx11$\,\el): with the correct read noise
the background model fits them and GMX would run there (pairs 0.15 $\to$ 0.97), but its footprint energy collects
$\sqrt{N_\mathrm{pix}}\times3.5$\,\el{} and the 59.5\,keV line widens by $\sim40\,\%$; choosing proper thresholds and
energies for noisy images is an open item.

\paragraph{Found sound.} The Poisson$\otimes$Gaussian pixel model and its variance $\lambda+s^2$; the matched filter;
the exact Poisson box check; charge conservation and unbiased energies (also at 1\,\el/pix and at every width); one
dark-charge subtraction everywhere; units; the satellite limit, the width-ratio rule and the CTI growth.

\paragraph{Documented, not changed.}
\begin{itemize}
\item Detection fakes depend on the dark charge: none in 4\,Mpix at $\lambda\le0.05$, about one 16--27-\el{} cluster
per 3--4\,Mpix at $\lambda=1$--2\,\el/pix.
\item Footprint noise: $\sigma_E^2=\sum_p\tilde r_p^2V_p$ over the $4\sigma$ ellipse; K$\alpha$ ($\sigma=1.5$\,px): 2.2\,\el{}
at $\lambda=0.01$, 11\,\el{} at $\lambda=1$, comparable with the Fano width 13.7\,\el. A threshold sum has less noise
but loses the tails.
\item The ridge test merges equal clouds up to $d=3.3\,\sigma_\mathrm{tot}$ although they are bimodal above
$2\,\sigma_\mathrm{tot}$ ($\approx3.5$\,px back side, 1.4\,px front side); only \code{peakfrac} shows such merges.
\item The island BIC test uses a heuristic $\chi^2$ likelihood with an un-derived margin of 30, compares a global model
with an island refit (biased toward merging), and its round-X-ray hypothesis penalises elliptical X-rays (CTI, column
binning). To be recalibrated on injected tracks before changing it.
\item The blend cut removes deep (diffuse) events preferentially in crowded fields: fine for line positions, an
efficiency correction for rates.
\item In the WADERs-format files, \AGMX{} Energy is dark-subtracted (WADERs includes $N_\mathrm{pix}\lambda$), and
\code{STD\_X/Y}, \code{DX/DY}, \code{npix} are footprint sizes: WADERs cuts on them do not transfer (\code{fSTD} does).
\item Test harness: the simulated muons deposited 250\,\el/px (0.94\,keV/px), four times below a minimum-ionising
particle.
\end{itemize}

%======================================================================================================
\section{Few-electron events at the lowest dark current (26 October)}\label{sec:lowdark}

For rare-event searches the decisive regime is a dark current of $10^{-4}$--$10^{-3}$\,\el/pix. Three problems limited
\AGMX{} there: the $\lambda$ floor of $10^{-3}$; the analytic threshold, which counts one trial per pixel although at
such dark currents an image has only a handful of candidate groups (0.001\,\el/pix: 0.000 fakes per image instead of
the wanted 0.14, and 3-\el{} efficiency 0.77 instead of 0.93); and the fitted $\lambda$ map, which fluctuates far
below and above the truth when a fit cell holds only a few dark electrons. The detection background is now the dark
charge implied by the measured rate of \emph{isolated} single-electron pixels (dark electrons are isolated; real
events come with neighbours), smoothed over $\pm30$ rows, with the rounding miscounts of the read noise; \AGMX's
per-event finder uses it conservatively (never below the fitted map), the candidate catalogue
(section~\ref{sec:fecand}) exactly.

\begin{table}[H]\centering\small
\begin{tabular}{lcccc}
\toprule
dark charge [\el/pix] & 3\,\el{} found (before $\to$ now) & 4\,\el{} found & fakes per 1.7\,Mpix (before $\to$ now) \\
\midrule
0.0001 & 0.91 $\to$ \textbf{0.96} & 1.00 & 0.00 $\to$ 0.30 \\
0.0005 & 0.90 $\to$ \textbf{0.93} & 0.99 & 0.00 $\to$ 0.37 \\
0.001 & 0.84 $\to$ 0.85 & 1.00 & 0.07 $\to$ 0.07 \\
0.003 & 0.39 $\to$ 0.40 & 0.98 $\to$ 0.96 & 0.30 $\to$ \textbf{0.07} \\
0.01 & 0.04 & 0.37 $\to$ \textbf{0.53} & 0.43 $\to$ \textbf{0.23} \\
\bottomrule
\end{tabular}
\caption*{Full \AGMX{} on simulated images (0.15\,\el{} read noise): 1800 events of 2--5\,\el{} (Gaussian
$\sigma$ 0.2--0.6\,px) and 30 blank images per dark charge.}
\end{table}

What a per-event decision cannot do at the lowest dark current: the remaining fakes there are mostly single pixels
with 2\,\el. Dark charge and rounding miscounts produce about 0.5 of them per 1.7-Mpix image at
$5\times10^{-4}$\,\el/pix ($\lambda^2/2$ plus $\lambda\,p$), indistinguishable from a real 2-\el{} event; two
neighbouring single electrons (about 14 per image at $10^{-3}$\,\el/pix) likewise. Rare-event searches therefore do
not decide per event: they keep every candidate and compare the counts with the expected background, which is what the
candidate catalogue provides.

%======================================================================================================
\section{The few-electron candidate catalogue (26 October)}\label{sec:fecand}

A separate stage after \AGMX{} (\code{clustering\_tests/fe\_candidates.py}), in the style of the low-background
experiments (DAMIC-M, SENSEI). For electron-counting images (read noise $<0.35$\,\el) \code{run\_AGMX} and
\code{run\_WADERs --clusterer agmx} write \code{fe\_candidates/<image>\_ext<n>\_fe\_candidates.csv} and \code{.json}:
\begin{itemize}
\item \textbf{candidates}: every group of touching pixels with 2--20 counted electrons ($\le9$ pixels): position,
charge, pixels, local dark charge and read noise, the statistic $T$, $n_\mathrm{exp}$ = the expected number of
dark-charge groups per image at least as unlikely in the clean area (simulated blanks; analytic tail
$N_\mathrm{clean}e^{-T}$ beyond them) and $p_\mathrm{image}=1-e^{-n_\mathrm{exp}}$;
\item \textbf{mask flags} (nothing is removed): halo within $5\sqrt{E/\mathrm{keV}}$\,px (5--60\,px) of an \AGMX{}
cluster, charge-transfer trails along its row and column ($E>1$\,keV, 50\,px), hot columns (also from an excess
single-electron rate), the edge (5\,px), overlap with an \AGMX{} cluster; \code{clean} = no flag;
\item \textbf{summary}: the isolated single-electron rate, the clean area, and per class (charge, pixels) the observed
clean candidates and the number expected from dark charge alone (blanks simulated from the image's own detection
background).
\end{itemize}

\begin{table}[H]\centering\small
\begin{tabular}{lccc}
\toprule
dark charge [\el/pix] & 2\,\el{} in 1\,px & 2\,\el{} in 2\,px & 3\,\el{} in 3\,px \\
\midrule
0.0005 & 5 / 5.2 & 58 / 48.2 & 0 / 0.1 \\
0.001 & 14 / 12.7 & 120 / 114.8 & 0 / 0.7 \\
0.003 & 89 / 93.4 & 731 / 734.0 & 11 / 13.4 \\
\bottomrule
\end{tabular}
\caption*{Calibration: observed / expected clean candidates in 10 blank images (1.7\,Mpix each) per dark charge.
Injected 3-\el{} events are listed in 99\,\% of cases (median $n_\mathrm{exp}=0.008$ per image); candidates next to a
60-keV deposit carry the halo flag.}
\end{table}
